\documentclass[12pt,a4paper]{article}
\usepackage[a4paper,margin=18mm]{geometry}
\usepackage[T2A]{fontenc}
\usepackage[utf8]{inputenc}
\usepackage[russian]{babel}
\usepackage{amsmath,amssymb,mathtools}
\usepackage{array,booktabs,longtable}
\usepackage{xcolor}
\usepackage{hyperref}
\usepackage{tikz}

\hypersetup{
  colorlinks=true,
  linkcolor=blue!55!black,
  urlcolor=blue!55!black
}
\setlength{\parindent}{0pt}
\setlength{\parskip}{5pt}
\setlength{\emergencystretch}{2em}
\renewcommand{\arraystretch}{1.18}

\begin{document}

\begin{titlepage}
\centering
\vspace*{0.18\textheight}
{\LARGE\bfseries Ревью домашней работы\par}
\vspace{18mm}
{\large Автор: Лёвин Артём Александрович\par}
\vspace{6mm}
{\large 12 сентября 2026 года\par}
\vfill
\begin{tikzpicture}
  \draw[line width=0.7pt] (0,0) .. controls (2.2,0.55) and (5.3,-0.55) .. (7.8,0);
\end{tikzpicture}
\end{titlepage}

\newpage
\section*{Сводная таблица}
\small
\begin{longtable}{|p{0.8cm}|p{2.7cm}|p{4.5cm}|p{2.5cm}|p{2.6cm}|}
\hline
\textbf{№} & \textbf{Тема} & \textbf{Суть ошибки} & \textbf{Ваш ответ} & \textbf{Правильный ответ} \\
\hline
\hyperref[task:2]{2} & Раскрытие скобок & В промежуточной строке $-2(x+5)$ раскрыто как $-2x+10$. Конечный ответ исправлен, но записанный переход неверен. & $4x-34$ & $4x-34$ \\
\hline
\hyperref[task:3]{3} & Минус перед скобками & При раскрытии $-(3x-7)$ знак у $7$ не изменён. Конечный ответ исправлен, но из записанной строки он не следует. & $x+9$ & $x+9$ \\
\hline
\hyperref[task:4]{4} & Раскрытие скобок & В обеих скобках внешний множитель применён не ко всем слагаемым с нужным коэффициентом. & $9x-6y$ & $11x-15y$ \\
\hline
\hyperref[task:5]{5} & Упрощение и подстановка & В $-2(x-4y)$ получено $-2x+6y$ вместо $-2x+8y$, поэтому неверны и упрощение, и значение. & $-x+14y$, затем $16$ & $-x+16y$, затем $18$ \\
\hline
\hyperref[task:7]{7} & Приведение подобных & Скобки раскрыты верно, но $15-4$ вычислено как $9$ вместо $11$. & $-2x+9$ & $-2x+11$ \\
\hline
\hyperref[task:9]{9} & Минус перед скобками & В $-(2x-3y)$ знак перед $3y$ изменён неверно: должно быть $+3y$. & $3x-8y$ & $3x-2y$ \\
\hline
\hyperref[task:10]{10} & Упрощение и подстановка & В $-3(x+2y)$ записано $-3x-5y$ вместо $-3x-6y$; окончательное числовое значение не доведено до ответа. & $-x-7y$; далее запись не совсем однозначна & $-x-8y$, затем $17$ \\
\hline
\end{longtable}
\normalsize

\newpage
\thispagestyle{empty}
\vspace*{\fill}
\begin{center}
{\LARGE\bfseries Разбор ошибок}
\end{center}
\vspace*{\fill}

\newpage
\phantomsection\label{task:2}
\section*{Задание 2}

\textbf{Текст задания.}
Раскройте скобки и упростите:
\[
6(x-4)-2(x+5).
\]

\textbf{Ваше решение.}
В записи видно:
\[
6x-24-2x+10,
\]
после чего промежуточный результат зачёркнут и справа записан итог
\[
4x-34.
\]

\textbf{Ваш ответ.}
\[
4x-34.
\]

\textcolor{red}{Ошибка:} неверный знак при раскрытии второй скобки.

\textbf{Где именно возникает ошибка.}
Последний корректный фрагмент --- исходное выражение
\[
6(x-4)-2(x+5).
\]
Первая скобка раскрыта верно:
\[
6(x-4)=6x-24.
\]
Первый достоверно неверный переход появляется во второй скобке:
\[
-2(x+5)=-2x+10.
\]
Здесь произведение $-2\cdot 5$ записано со знаком «плюс».

\textbf{Почему этот шаг неверен.}
По распределительному закону множитель перед скобками умножается на \emph{каждое} слагаемое внутри скобок:
\[
-2(x+5)=(-2)\cdot x+(-2)\cdot 5.
\]
Первое произведение равно $-2x$, а второе равно $-10$, потому что отрицательное число, умноженное на положительное, даёт отрицательное число. Поэтому
\[
-2(x+5)=-2x-10.
\]
Записанный в конце ответ $4x-34$ верный, но из строки с $+10$ он не получается. Значит, конечное исправление сделано правильно, а ход решения необходимо привести в соответствие с ответом.

\textcolor{blue!60!black}{Ключевая идея:} при раскрытии скобок сначала отдельно умножайте внешний множитель на каждый член скобки и только затем приводите подобные слагаемые.

\textbf{Исправление от последнего правильного шага.}
Нужно заменить только неверное раскрытие второй скобки:
\[
6(x-4)-2(x+5)=6x-24-2x-10.
\]
Теперь можно привести подобные слагаемые:
\[
6x-2x=4x,
\qquad
-24-10=-34.
\]
Отсюда
\[
4x-34.
\]

\textbf{Полное правильное решение.}
\[
\begin{aligned}
6(x-4)-2(x+5)
&=6x-24-2x-10\\
&=(6x-2x)+(-24-10)\\
&=4x-34.
\end{aligned}
\]

\textbf{Проверка.}
Возьмём $x=0$. Исходное выражение:
\[
6(0-4)-2(0+5)=-24-10=-34.
\]
Полученное выражение:
\[
4\cdot 0-34=-34.
\]
Значения совпадают.

\textcolor{teal!60!black}{Итог:} конечный ответ у Вас исправлен верно, но в раскрытии второй скобки нужно заменить $+10$ на $-10$.

\textbf{Задача на закрепление.}
Упростите выражение:
\[
5(x-3)-3(x+4).
\]

\newpage
\phantomsection\label{task:3}
\section*{Задание 3}

\textbf{Текст задания.}
Раскройте скобки и упростите:
\[
-(3x-7)+2(2x+1).
\]

\textbf{Ваше решение.}
В записи получено
\[
-3x-7+4x+2,
\]
после чего промежуточный ответ исправлен и записано
\[
x+9.
\]

\textbf{Ваш ответ.}
\[
x+9.
\]

\textcolor{red}{Ошибка:} неверно раскрыта скобка, перед которой стоит минус.

\textbf{Где именно возникает ошибка.}
Первый неверный переход:
\[
-(3x-7)=-3x-7.
\]
Коэффициент перед скобкой здесь равен $-1$. Он должен умножаться на оба слагаемых внутри скобки.

\textbf{Почему этот шаг неверен.}
Запись
\[
-(3x-7)
\]
означает
\[
(-1)(3x-7).
\]
Поэтому
\[
(-1)\cdot 3x=-3x,
\qquad
(-1)\cdot(-7)=+7.
\]
Знак у второго слагаемого обязан измениться с «минус» на «плюс». Следовательно,
\[
-(3x-7)=-3x+7.
\]
Ваша конечная запись $x+9$ совпадает с правильным ответом, но она не следует из промежуточной строки $-3x-7+4x+2$: из этой строки получилось бы $x-5$. Поэтому здесь важно исправить именно ход преобразований.

\textcolor{blue!60!black}{Ключевая идея:} минус перед скобкой означает умножение каждого слагаемого внутри скобки на $-1$, поэтому знаки всех слагаемых меняются на противоположные.

\textbf{Исправление от последнего правильного шага.}
\[
-(3x-7)+2(2x+1)=-3x+7+4x+2.
\]
Теперь приводим подобные:
\[
-3x+4x=x,
\qquad
7+2=9.
\]
Получаем
\[
x+9.
\]

\textbf{Полное правильное решение.}
\[
\begin{aligned}
-(3x-7)+2(2x+1)
&=-3x+7+4x+2\\
&=(-3x+4x)+(7+2)\\
&=x+9.
\end{aligned}
\]

\textbf{Проверка.}
При $x=0$ исходное выражение равно
\[
-(-7)+2\cdot 1=7+2=9.
\]
Упрощённое выражение даёт
\[
0+9=9.
\]
Совпадение подтверждает преобразование.

\textcolor{teal!60!black}{Итог:} ответ $x+9$ верный, но в строке раскрытия скобок необходимо писать $-3x+7$, а не $-3x-7$.

\textbf{Задача на закрепление.}
Упростите выражение:
\[
-(4x-9)+3(x+2).
\]

\newpage
\phantomsection\label{task:4}
\section*{Задание 4}

\textbf{Текст задания.}
Упростите выражение:
\[
4(2x-3y)-3(y-x).
\]

\textbf{Ваше решение.}
\[
8x-3y-3y+x=-6y+9x.
\]

\textbf{Ваш ответ.}
\[
9x-6y.
\]

\textcolor{red}{Ошибка:} внешние множители применены не ко всем слагаемым в скобках с нужными коэффициентами.

\textbf{Где именно возникает ошибка.}
Первый неверный шаг появляется уже при раскрытии первой скобки:
\[
4(2x-3y)\rightarrow 8x-3y.
\]
Часть $4\cdot 2x=8x$ выполнена верно, но второе слагаемое должно также умножаться на $4$:
\[
4\cdot(-3y)=-12y.
\]
Кроме того, во второй скобке допущена ещё одна такая же ошибка:
\[
-3(y-x)\rightarrow -3y+x,
\]
хотя
\[
(-3)\cdot(-x)=+3x.
\]

\textbf{Почему это неверно.}
Распределительный закон имеет вид
\[
a(b+c)=ab+ac.
\]
То есть внешний множитель нельзя применить только к первому слагаемому и оставить второе без изменения. В данном выражении нужно выполнить четыре произведения:
\[
4\cdot 2x,
\quad
4\cdot(-3y),
\quad
(-3)\cdot y,
\quad
(-3)\cdot(-x).
\]
Они равны соответственно
\[
8x,
\quad
-12y,
\quad
-3y,
\quad
3x.
\]

\textcolor{blue!60!black}{Ключевая идея:} перед приведением подобных слагаемых сначала полностью раскройте каждую скобку, умножив внешний коэффициент на каждый её член.

\textbf{Исправление от последнего правильного шага.}
Исходное выражение переписываем с корректно раскрытыми скобками:
\[
4(2x-3y)-3(y-x)=8x-12y-3y+3x.
\]
Теперь объединяем отдельно слагаемые с $x$ и отдельно с $y$:
\[
8x+3x=11x,
\qquad
-12y-3y=-15y.
\]
Поэтому
\[
11x-15y.
\]

\textbf{Полное правильное решение.}
\[
\begin{aligned}
4(2x-3y)-3(y-x)
&=8x-12y-3y+3x\\
&=(8x+3x)+(-12y-3y)\\
&=11x-15y.
\end{aligned}
\]

\textbf{Проверка.}
При $x=1$, $y=1$ исходное выражение равно
\[
4(2-3)-3(1-1)=4(-1)-0=-4.
\]
Правильное упрощённое выражение даёт
\[
11-15=-4.
\]
Ваш результат $9x-6y$ при тех же значениях дал бы $3$, поэтому он не равносилен исходному выражению.

\textcolor{teal!60!black}{Итог:} правильно упрощённое выражение равно $11x-15y$; ключевая ошибка --- неполное применение множителей к членам скобок.

\textbf{Задача на закрепление.}
Упростите выражение:
\[
5(2x-4y)-2(y-3x).
\]

\newpage
\phantomsection\label{task:5}
\section*{Задание 5}

\textbf{Текст задания.}
Сначала упростите выражение, затем найдите его значение при
\[
x=-2,\qquad y=1:
\]
\[
3(2x+y)-2(x-4y)+5(y-x).
\]

\textbf{Ваше решение.}
\[
6x+3y-2x+6y+5y-5x=-x+14y,
\]
после чего при подстановке получено
\[
2+14=16.
\]

\textbf{Ваш ответ.}
\[
16.
\]

\textcolor{red}{Ошибка:} неверно умножено $-2$ на $-4y$.

\textbf{Где именно возникает ошибка.}
Первая скобка раскрыта верно:
\[
3(2x+y)=6x+3y.
\]
Первый неверный шаг находится во второй скобке:
\[
-2(x-4y)\rightarrow -2x+6y.
\]
Правильное произведение равно
\[
(-2)\cdot(-4y)=+8y,
\]
поэтому должно быть
\[
-2x+8y.
\]

\textbf{Почему этот шаг неверен.}
При раскрытии скобок нужно умножить $-2$ отдельно на $x$ и на $-4y$. Знак второго произведения положительный, потому что перемножаются два отрицательных множителя, а коэффициент равен $2\cdot 4=8$. Ошибка в коэффициенте $6$ уменьшила сумму коэффициентов при $y$ на $2$, поэтому дальше получилось $14y$ вместо $16y$. Именно эта ошибка затем дала числовой ответ $16$ вместо $18$.

\textcolor{blue!60!black}{Ключевая идея:} при раскрытии скобок контролируйте отдельно знак и числовой коэффициент каждого произведения; после упрощения только затем подставляйте значения переменных.

\textbf{Исправление от последнего правильного шага.}
Заменяем только неверный фрагмент:
\[
3(2x+y)-2(x-4y)+5(y-x)
=6x+3y-2x+8y+5y-5x.
\]
Собираем слагаемые с $x$:
\[
6x-2x-5x=-x.
\]
Собираем слагаемые с $y$:
\[
3y+8y+5y=16y.
\]
Значит,
\[
-x+16y.
\]
Теперь подставляем $x=-2$, $y=1$:
\[
-(-2)+16\cdot 1=2+16=18.
\]

\textbf{Полное правильное решение.}
\[
\begin{aligned}
3(2x+y)-2(x-4y)+5(y-x)
&=6x+3y-2x+8y+5y-5x\\
&=(6x-2x-5x)+(3y+8y+5y)\\
&=-x+16y.
\end{aligned}
\]
При $x=-2$, $y=1$:
\[
-x+16y=-(-2)+16\cdot 1=18.
\]

\textbf{Проверка.}
Проверим значение непосредственно по исходному выражению:
\[
3(2\cdot(-2)+1)-2((-2)-4\cdot 1)+5(1-(-2)).
\]
Вычисляем по частям:
\[
3(-3)-2(-6)+5\cdot 3=-9+12+15=18.
\]
Получено то же значение.

\textcolor{teal!60!black}{Итог:} правильное упрощение --- $-x+16y$, а правильное значение при заданных $x$ и $y$ равно $18$.

\textbf{Задача на закрепление.}
Сначала упростите выражение, затем найдите его значение при $a=-1$, $b=2$:
\[
2(3a+b)-3(a-2b)+4(b-a).
\]

\newpage
\phantomsection\label{task:7}
\section*{Задание 7}

\textbf{Текст задания.}
Раскройте скобки и упростите:
\[
-3(2x-5)+4(x-1).
\]

\textbf{Ваше решение.}
\[
-6x+15+4x-4=-2x+9.
\]

\textbf{Ваш ответ.}
\[
-2x+9.
\]

\textcolor{red}{Ошибка:} неверно вычислена разность свободных членов.

\textbf{Где именно возникает ошибка.}
Раскрытие обеих скобок выполнено правильно:
\[
-3(2x-5)=-6x+15,
\qquad
4(x-1)=4x-4.
\]
После этого строка
\[
-6x+15+4x-4
\]
полностью корректна. Первый неверный шаг появляется при приведении числовых слагаемых:
\[
15-4\rightarrow 9.
\]
На самом деле
\[
15-4=11.
\]

\textbf{Почему этот шаг неверен.}
Подобные слагаемые с $x$ и свободные числа нужно объединять независимо:
\[
-6x+4x=-2x,
\]
а свободная часть равна
\[
15+(-4)=11.
\]
Знак перед $4$ уже учтён в записи $-4$, поэтому здесь выполняется обычное вычитание четырёх из пятнадцати.

\textcolor{blue!60!black}{Ключевая идея:} после раскрытия скобок группируйте отдельно переменные и отдельно числа; это снижает риск смешать коэффициенты и свободные члены.

\textbf{Исправление от последнего правильного шага.}
Начинаем с уже верной строки:
\[
-6x+15+4x-4.
\]
Группируем:
\[
(-6x+4x)+(15-4)=-2x+11.
\]

\textbf{Полное правильное решение.}
\[
\begin{aligned}
-3(2x-5)+4(x-1)
&=-6x+15+4x-4\\
&=(-6x+4x)+(15-4)\\
&=-2x+11.
\end{aligned}
\]

\textbf{Проверка.}
При $x=0$ исходное выражение равно
\[
-3(-5)+4(-1)=15-4=11.
\]
Упрощённое выражение даёт
\[
-2\cdot 0+11=11.
\]

\textcolor{teal!60!black}{Итог:} раскрытие скобок сделано верно; исправить нужно только арифметику $15-4=11$.

\textbf{Задача на закрепление.}
Упростите выражение:
\[
-4(2x-3)+3(x-5).
\]

\newpage
\phantomsection\label{task:9}
\section*{Задание 9}

\textbf{Текст задания.}
Упростите выражение:
\[
-(2x-3y)-2(4y-x)+3(x+y).
\]

\textbf{Ваше решение.}
\[
-2x-3y-8y+2x+3x+3y=3x-8y.
\]

\textbf{Ваш ответ.}
\[
3x-8y.
\]

\textcolor{red}{Ошибка:} неверно изменён знак у $-3y$ при раскрытии первой скобки.

\textbf{Где именно возникает ошибка.}
Первый неверный переход:
\[
-(2x-3y)\rightarrow -2x-3y.
\]
Перед скобкой стоит коэффициент $-1$, поэтому оба слагаемых должны изменить знак:
\[
-(2x-3y)=-2x+3y.
\]
Остальные скобки в Вашей записи раскрыты правильно:
\[
-2(4y-x)=-8y+2x,
\qquad
3(x+y)=3x+3y.
\]

\textbf{Почему этот шаг неверен.}
Выражение $-(2x-3y)$ означает умножение всей скобки на $-1$:
\[
(-1)\cdot 2x=-2x,
\qquad
(-1)\cdot(-3y)=+3y.
\]
Второе слагаемое было отрицательным внутри скобки, поэтому после умножения на $-1$ оно становится положительным. Из-за сохранённого минуса коэффициент при $y$ оказался занижен на $6$: вместо суммы $3-8+3=-2$ в Вашей строке получается $-3-8+3=-8$.

\textcolor{blue!60!black}{Ключевая идея:} если перед скобкой стоит только знак «минус», воспринимайте его как множитель $-1$ и меняйте знак каждого слагаемого внутри скобки.

\textbf{Исправление от последнего правильного шага.}
Корректно раскрываем первую скобку и сохраняем верные части остальных:
\[
-2x+3y-8y+2x+3x+3y.
\]
Теперь приводим подобные слагаемые:
\[
-2x+2x+3x=3x,
\]
\[
3y-8y+3y=-2y.
\]
Следовательно,
\[
3x-2y.
\]

\textbf{Полное правильное решение.}
\[
\begin{aligned}
-(2x-3y)-2(4y-x)+3(x+y)
&=-2x+3y-8y+2x+3x+3y\\
&=(-2x+2x+3x)+(3y-8y+3y)\\
&=3x-2y.
\end{aligned}
\]

\textbf{Проверка.}
При $x=1$, $y=1$ исходное выражение равно
\[
-(2-3)-2(4-1)+3(1+1)=1-6+6=1.
\]
Упрощённое выражение даёт
\[
3\cdot 1-2\cdot 1=1.
\]

\textcolor{teal!60!black}{Итог:} правильный результат --- $3x-2y$; ошибка возникла только при смене знаков в первой скобке.

\textbf{Задача на закрепление.}
Упростите выражение:
\[
-(3x-4y)-2(2y-x)+2(x+y).
\]

\newpage
\phantomsection\label{task:10}
\section*{Задание 10}

\textbf{Текст задания.}
Сначала упростите выражение, затем найдите его значение при
\[
x=-1,\qquad y=-2:
\]
\[
4(2x-y)-3(x+2y)-2(3x-y).
\]

\textbf{Ваше решение.}
В основной строке записано
\[
8x-4y-3x-5y-6x+2y=-x-7y.
\]
После этого есть ещё одна короткая запись подстановки; не совсем понятно, что именно записано после знака равенства, поэтому окончательный числовой ответ надёжно восстановить нельзя. Для проверки это не мешает: первая математическая ошибка видна раньше.

\textbf{Ваш ответ.}
Упрощённое выражение записано как
\[
-x-7y,
\]
окончательный числовой ответ не указан однозначно.

\textcolor{red}{Ошибка:} во второй скобке неверно вычислен коэффициент при $y$.

\textbf{Где именно возникает ошибка.}
Первая скобка раскрыта верно:
\[
4(2x-y)=8x-4y.
\]
Первый неверный переход:
\[
-3(x+2y)\rightarrow -3x-5y.
\]
Нужно умножить $-3$ на $2y$:
\[
(-3)\cdot 2y=-6y,
\]
поэтому правильный результат раскрытия этой скобки:
\[
-3x-6y.
\]
Третья скобка раскрыта верно:
\[
-2(3x-y)=-6x+2y.
\]

\textbf{Почему этот шаг неверен.}
Коэффициент $2$ при $y$ является частью слагаемого $2y$. При умножении на $-3$ числовые коэффициенты перемножаются:
\[
(-3)\cdot 2=-6.
\]
Замена $-6y$ на $-5y$ уменьшает модуль коэффициента на единицу. Затем при сложении коэффициентов при $y$ получается $-7$ вместо $-8$, поэтому и дальнейшая подстановка уже не может дать правильный ответ.

\textcolor{blue!60!black}{Ключевая идея:} в произведении вида $a(bx+cy)$ каждый коэффициент внутри скобки умножается на $a$ полностью: отдельно контролируйте знак и произведение числовых коэффициентов.

\textbf{Исправление от последнего правильного шага.}
Заменяем $-5y$ на $-6y$:
\[
8x-4y-3x-6y-6x+2y.
\]
Собираем слагаемые с $x$:
\[
8x-3x-6x=-x.
\]
Собираем слагаемые с $y$:
\[
-4y-6y+2y=-8y.
\]
Получаем
\[
-x-8y.
\]
Теперь подставляем $x=-1$, $y=-2$:
\[
-(-1)-8(-2)=1+16=17.
\]

\textbf{Полное правильное решение.}
\[
\begin{aligned}
4(2x-y)-3(x+2y)-2(3x-y)
&=8x-4y-3x-6y-6x+2y\\
&=(8x-3x-6x)+(-4y-6y+2y)\\
&=-x-8y.
\end{aligned}
\]
При $x=-1$, $y=-2$:
\[
-x-8y=-(-1)-8(-2)=17.
\]

\textbf{Проверка.}
Подставим значения непосредственно в исходное выражение:
\[
4(2(-1)-(-2))-3((-1)+2(-2))-2(3(-1)-(-2)).
\]
Вычисляем скобки:
\[
2(-1)-(-2)=0,
\qquad
-1-4=-5,
\qquad
-3+2=-1.
\]
Тогда
\[
4\cdot 0-3(-5)-2(-1)=0+15+2=17.
\]
Результат совпадает.

\textcolor{teal!60!black}{Итог:} правильное упрощение --- $-x-8y$, а при $x=-1$, $y=-2$ значение выражения равно $17$.

\textbf{Задача на закрепление.}
Сначала упростите выражение, затем найдите его значение при $x=2$, $y=-1$:
\[
3(2x-y)-4(x+3y)-2(2x-y).
\]

\vfill
\begin{center}
\textit{Лёвин Артём Александрович эксклюзивно для Кирилла Зиновьева}
\end{center}

\end{document}
